% !TEX program = xelatex
\documentclass[11pt,a4paper]{article}

\usepackage{fontspec}
\setmainfont{Arial}
\usepackage[english]{babel}
\usepackage{microtype}
\usepackage[a4paper,margin=2.5cm]{geometry}
\usepackage{amsmath,amssymb}
\usepackage{booktabs}
\usepackage{tabularx}
\usepackage{array}
\usepackage{enumitem}
\usepackage{xcolor}
\usepackage{titlesec}
\usepackage[hidelinks,breaklinks]{hyperref}

\setlist{itemsep=2pt,topsep=4pt}
\setlength{\parskip}{0.5em}
\setlength{\parindent}{0pt}
\titleformat{\section}{\normalfont\large\bfseries}{}{0pt}{}
\titlespacing*{\section}{0pt}{1.4em}{0.5em}

\newcommand{\gpuzip}{GPUZIP}

\begin{document}

\begin{center}
{\Large\bfseries Accurate FP64 Seismic Simulation on Modern NVIDIA
Architectures: Co-designing Mixed Precision, Compressed Checkpointing, and
Prefetch for PDE-Based Reverse Time Migration}
\end{center}

\vspace{1em}

\noindent\textbf{Principal Investigator:} Prof.\ Sandro Rigo\\
\textbf{Title/Role:} Associate Professor, Institute of Computing\\
\textbf{University:} Universidade Estadual de Campinas\\
\textbf{Collaborator 1:} Thiago Maltempi\,---\,Graduate Researcher (PhD
Candidate), Institute of Computing, Universidade Estadual de Campinas\\
\textbf{Collaborator 2:} Prof.\ Hervé Yviquel\,---\,Assistant Professor,
Institute of Computing, Universidade Estadual de Campinas

% ============================================================================
\section*{Abstract}

Scientific simulations such as Reverse Time Migration (RTM) solve partial
differential equations (PDEs) demanding FP64/FP32 accuracy, yet modern NVIDIA
GPUs increasingly optimize for low-precision AI workloads (FP8/BF16) and expose
new memory hierarchies, tensor cores, and dense linear-algebra libraries with
mixed-precision extensions. Our prior work, \textbf{\gpuzip}, accelerated RTM on
V100 systems by combining GPU-side lossy compression (NVIDIA Bitcomp via nvCOMP)
with checkpoint prefetching, cutting host--GPU transfer volume by up to
$80\times$ and delivering up to $5.1\times$ end-to-end speedup. This project
extends \gpuzip{} to modern NVIDIA hardware over a 12-month programme in two
stages. In the first six months, using H100 80\,GB cloud GPUs, we will:
(1)~port Awave-3D and \gpuzip{} to the NVIDIA HPC SDK; and (2)~re-evaluate
compressed checkpointing and prefetch on Hopper memory, introducing
\textbf{differential checkpointing}\,---\,a hybrid scheme that intermixes full
wavefield snapshots with compressed inter-timestep deltas to reduce both
transfer volume and recomputation cost. In the second six months, using an RTX
PRO 6000 Blackwell GPU, we will: (3)~introduce mixed-precision wavefield updates
via cuBLAS mixed-precision GEMM and cuSOLVER iterative refinement; and (4)~map
where Ozaki-scheme FP64 emulation via low-precision tensor cores,
cuFFT-accelerated spectral operators, and the precision$\times$compression
interaction improve a memory-bound stencil solver on AI-first hardware.
Outcomes include an open-source \gpuzip{} v3.0 release and a peer-reviewed
publication.

\textbf{Keywords:} high-performance computing; reverse time migration; mixed
precision; lossy compression; checkpointing

% ============================================================================
\section*{NVIDIA Platforms}

\textbf{HPC SDK} (nvc++/nvfortran, CUDA-X math libraries, Nsight
Systems/Compute). \emph{Experience:} team uses CUDA and Nsight Systems
routinely. \emph{Adoption:} migrate Awave-3D build from clang to nvc++; validate
numerical parity against existing baselines. Nsight Compute provides
kernel-level roofline analysis to identify memory-bandwidth-bound vs.\
compute-bound kernels and objectively scope tensor-core opportunities.

\textbf{nvCOMP / Bitcomp} (GPU-side lossy compressor). \emph{Experience:}
integrated and validated in \gpuzip{} v2.0; Bitcomp delivered the best
accuracy-compression balance on V100. \emph{Adoption:} re-tune error bounds for
full-checkpoint and delta-checkpoint paths; re-benchmark throughput against HBM3
bandwidth on H100.

\textbf{cuBLAS (mixed-precision extensions) and cuSOLVER (iterative-refinement
solver).} Both ship inside the HPC SDK. \emph{Experience:} team has used BLAS
and LAPACK routines in prior HPC work. \emph{Adoption:} apply cuBLAS
mixed-precision GEMM (FP32 compute with FP64 accumulation) and cuSOLVER
iterative refinement to dense sub-kernels in the RTM pipeline. Evaluate
Ozaki-scheme FP64 emulation via low-precision tensor-core units on the RTX PRO
6000 Blackwell\,---\,a GPU whose negligible native FP64 makes it the ideal
platform to characterise whether tensor-core emulation can substitute for
hardware FP64 in scientific workloads.

\textbf{cuFFT.} \emph{Experience:} used in prior spectral operator work.
\emph{Adoption:} benchmark cuFFT-based pseudo-spectral derivative operators as
an alternative to FDM stencils for selected wavefield updates; characterise the
memory-bandwidth vs.\ compute trade-off relative to the standard
finite-difference path on H100 HBM3.

\textbf{H100 80\,GB SXM5 (cloud) and RTX PRO 6000 Blackwell (hardware).} H100 is
used for Phases 1--2 (FP64-heavy, multi-GPU HPC requiring HBM3 and NVLink). RTX
PRO 6000 is used for Phases 3--4 (mixed-precision and Ozaki-scheme study, local
development, and continuation beyond the cloud expiry window). The RTX PRO 6000
Blackwell's 5th-generation tensor cores with FP8/FP4 support and 96\,GB GDDR7
ECC memory make it a scientifically distinct platform for Ozaki-scheme
emulation. Recent work has demonstrated up to $2.3\times$ speedup over native
FP64 DGEMM on this specific GPU.

% ============================================================================
\section*{Dataset and Model}

\begin{itemize}
  \item \textbf{Size:} \emph{$\approx$X\,TB total (confirm from REDU).}
        Checkpoint pools reached $\approx$330\,GB host RAM on V100; single
        Marmousi3D checkpoint $\approx$500\,MB.
  \item \textbf{Data type:} seismic velocity models and traces (numerical
        grids).
  \item \textbf{Source:} Salt (SEG/EAGE) and Marmousi3D are open; Large dataset
        is in-house.
  \item \textbf{Confidentiality:} non-confidential, no personal data, no NDA.
  \item \textbf{AI model parameters:} not applicable (numerical PDE solver).
  \item \textbf{Storage location:} Universidade Estadual de Campinas (UNICAMP),
        Campinas, Brazil; staged for upload to cloud.
\end{itemize}

% ============================================================================
\section*{Introduction}

Hardware is trending toward low precision for AI, but PDE-based scientific
simulation requires high precision to remain trustworthy. RTM, a seismic imaging
method that solves the wave equation by finite differences over thousands of
time steps, is a canonical example: geophysicists use the migrated image to
guide expensive drilling decisions, so numerical fidelity is non-negotiable. RTM
is also memory-intensive, requiring hundreds of gigabytes of wavefield state and
relying on checkpointing to trade memory against recomputation in its
forward/backward adjoint structure.

Our prior work addressed the dominant bottleneck (host--GPU checkpoint transfer
over PCIe) with \textbf{\gpuzip}~[1][2], which compresses checkpoints on the GPU
(Bitcomp) and proactively prefetches them, reaching up to $5.1\times$ speedup on
V100 while preserving image quality (SSIM $\approx 1$, near-zero relative
error). Three shifts now motivate this project. First, \textbf{the memory
bottleneck is changing}: Hopper-class coherent interconnects (HBM3, NVLink-C2C)
shrink the PCIe wall, making it an open empirical question whether compressed
checkpointing still pays and how the optimal cache size and prefetch schedule
shift. Second, \textbf{checkpoint strategies can be extended}: rather than
storing only full wavefield snapshots, a hybrid scheme can intermix full
checkpoints with compressed inter-timestep deltas (small difference fields) that
allow partial forward recomputation from the nearest full checkpoint without
restarting from scratch, reducing both storage volume and recomputation cost.
Third, \textbf{modern math libraries reshape the precision/performance trade
space}: cuBLAS exposes mixed-precision GEMM with FP64-accumulation paths,
cuSOLVER offers iterative-refinement solvers, and the Ozaki scheme enables
FP64-accurate DGEMM via low-precision tensor cores, an approach recently shown
to achieve up to $2.3\times$ speedup over native FP64 DGEMM on the RTX PRO 6000
Blackwell~[3]. To our knowledge, this is the first systematic co-design of
differential checkpointing, lossy compression, mixed-precision linear algebra,
and prefetch for an adjoint PDE solver on modern NVIDIA hardware.

% ============================================================================
\section*{Methods}

\textbf{Phase 1\,---\,HPC SDK port and validation (M1--M2, H100).} Build
Awave-3D and \gpuzip{} with nvc++, re-integrate nvCOMP/Bitcomp, and reproduce
V100 baselines on H100. Profile with Nsight Systems and Nsight Compute to
establish roofline positions for all kernels, confirming that FDM stencils are
memory-bandwidth-bound before any changes are attempted.

\textbf{Phase 2\,---\,Differential checkpointing and compression on Hopper
(M2--M5, H100).} Re-characterise the \gpuzip{} checkpoint pipeline on Hopper
(HBM3 bandwidth, larger L2, NVLink) and introduce \emph{differential
checkpointing}: a hybrid strategy that stores full compressed snapshots at
coarse intervals and compressed inter-timestep delta fields at finer
granularity. During the backward pass, deltas are accumulated from the nearest
full checkpoint, avoiding full forward recomputation for intermediate
timesteps. We characterise the accuracy impact of applying lossy compression
independently to both full checkpoints and deltas (PSNR, SSIM, mean relative
error), sweep GPU Checkpoint Cache size and prefetch schedule, and determine the
optimal full-to-delta checkpoint ratio across all three datasets and
checkpointing algorithms (Revolve, zCut, Uniform).

\textbf{Phase 3\,---\,Mixed-precision wavefield updates (M7--M10, RTX PRO
6000).} Introduce reduced-precision compute paths in the finite-difference
stencil (FP32 update with FP64 accumulation). For dense sub-kernels identified
in Phase 1, apply cuBLAS mixed-precision GEMM and cuSOLVER iterative
refinement; evaluate Ozaki-scheme FP64 emulation via low-precision tensor-core
units as an alternative. The RTX PRO 6000 Blackwell\,---\,with negligible native
FP64 but abundant FP8/FP4 tensor-core throughput\,---\,is the ideal platform for
this study: it forces the question of whether Ozaki-scheme emulation can fully
substitute for hardware FP64. Bound all errors jointly (compression $\times$
precision interaction) against the FP64 baseline using PSNR, SSIM, and mean
relative error.

\textbf{Phase 4\,---\,Spectral operators and final evaluation (M10--M12, RTX PRO
6000).} Benchmark cuFFT-based pseudo-spectral derivative operators as an
alternative to FDM stencils for selected wavefield updates. Integrate findings
from Phases 2--3 into a unified \gpuzip{} v3.0 release, covering differential
checkpointing, mixed-precision paths, and Ozaki-scheme emulation. Multi-node
scaling evaluation ($\leq$8 GPUs). The deliverable is a precise map of which
NVIDIA library and algorithm paths improve a memory-bound adjoint PDE solver and
which do not.

\textbf{Evaluation.} All experiments span three datasets (Large, Marmousi3D,
Salt) and three checkpointing algorithms (Revolve, zCut, Uniform). Kernel-level
profiling: Nsight Compute (roofline, memory-traffic breakdowns); timeline
profiling: Nsight Systems. Implementation: C/C++ with CUDA under the HPC SDK;
OpenMP Cluster for multi-node scaling.

% ============================================================================
\section*{Expected Results}

Timeline (12 months: M1--M6 on H100 cloud, M7--M12 on RTX PRO 6000):

\begin{center}
\small
\begin{tabularx}{\textwidth}{@{}l X l c@{}}
\toprule
\textbf{Month} & \textbf{Activity} & \textbf{Hardware} & \textbf{Phase} \\
\midrule
\textbf{M1--M2} & Port Awave-3D + \gpuzip{} to HPC SDK (nvc++); reproduce
  V100 baselines; Nsight Compute roofline profiling. & H100 (ramp-up) & 1 \\
\addlinespace
\textbf{M2--M5} & Differential checkpointing + lossy compression on Hopper;
  prefetch schedule optimisation; accuracy study (PSNR/SSIM/MRE). & H100
  (peak) & 2 \\
\addlinespace
\textbf{M5--M6} & Multi-node scaling ($\leq$8 GPUs); open-source \gpuzip{}
  v3.0-alpha release; reproducibility package (REDU). & H100 (8 GPUs) & 2 \\
\addlinespace
\textbf{M7--M10} & Mixed-precision stencil updates (cuBLAS/cuSOLVER/Ozaki);
  FP64 emulation on RTX PRO 6000; precision$\times$compression interaction. &
  RTX PRO 6000 & 3 \\
\addlinespace
\textbf{M10--M12} & cuFFT spectral-operator benchmark; unified \gpuzip{}
  v3.0 release; analysis, write-up, and paper submission. & RTX PRO 6000 &
  4 \\
\bottomrule
\end{tabularx}
\end{center}

\textbf{Projected outcomes:} open-source \gpuzip{} v3.0 targeting Hopper and
Blackwell NVIDIA architectures (MIT license, GitHub); peer-reviewed publication
(target: SBAC-PAD / Euro-Par / IJHPCA); public reproducibility package
(datasets, Nsight profiles, quality metrics) via UNICAMP REDU.

\textbf{Dependencies:} the project is not conditional on this award being
combined with another grant. Existing institutional support funds personnel and
local infrastructure; none duplicates the requested NVIDIA cloud compute or
hardware.

% ============================================================================
\section*{Project Support Details}

We request both cloud H100 80\,GB hours (Phases 1--2) and one RTX PRO 6000
Blackwell GPU (Phases 3--4). The H100 is required for FP64-heavy multi-GPU
experiments demanding Hopper's HBM3 bandwidth, FP64 Tensor Core throughput, and
NVLink fabric. The RTX PRO 6000 is requested for the mixed-precision and
Ozaki-scheme study in the second six months: its abundant FP8/FP4 tensor-core
throughput and 96\,GB GDDR7 ECC memory make it a scientifically distinct target
for FP64 emulation research, and it provides local compute beyond the cloud
expiry window.

\begin{center}
\small
\begin{tabularx}{\textwidth}{@{}X l l@{}}
\toprule
\textbf{Resource} & \textbf{Request} & \textbf{Cap} \\
\midrule
Cloud GPU hours (H100 80\,GB) & \textbf{$\sim$17,000 (Phases 1--2)} & 30,000 \\
Concurrent GPUs & \textbf{up to 8} & 8 \\
Cloud storage & \textbf{$\sim$4\,TB} & 32\,TB \\
Physical hardware & \textbf{1$\times$ RTX PRO 6000 Blackwell (Phases 3--4)} &
  $\leq$8 GPUs \\
\bottomrule
\end{tabularx}
\end{center}

\textbf{GPU-hour breakdown} (placeholders\,---\,replace with measured pilot
timings before submission):

\begin{itemize}
  \item \textbf{Phase 1 port + roofline profiling:} $\sim$2,000 GPU-hrs (H100)
  \item \textbf{Phase 2 differential checkpointing + compression study:}
        $\sim$13,000 GPU-hrs (H100, bulk)
  \item \textbf{Profiling overhead (Nsight Compute sweeps):} $\sim$2,000
        GPU-hrs (H100)
\end{itemize}

\textbf{Total $\approx$ 17,000 H100 hours}, well under the 30,000-hour cap.
Phases 3--4 run on the RTX PRO 6000 (local).

% ============================================================================
\section*{Cloud Readiness}

Development environment is Linux (Ubuntu); Awave-3D and \gpuzip{} already build
and run on Ubuntu, matching the cloud servers.

\begin{itemize}
  \item \textbf{Importing data:} pull datasets from UNICAMP REDU on day one.
        (Experienced.)
  \item \textbf{Installing dependencies:} HPC SDK (nvc++), nvCOMP, cuBLAS,
        cuSOLVER, cuFFT\,---\,all via NVIDIA NGC containers. (Experienced.)
  \item \textbf{Building/running the codebase:} existing CMake build scripts and
        CLI workflow. (Experienced.)
  \item \textbf{Backing up results:} Nsight profiles and checkpoint outputs
        synced off ephemeral instances regularly. (Experienced.)
  \item \textbf{CUDA and NVIDIA SDKs:} HPC SDK + CUDA-X + nvCOMP integration is
        core to the project. (Experienced.)
  \item \textbf{Containers:} containerised builds from NVIDIA NGC base images
        for full reproducibility. (Experienced.)
\end{itemize}

% ============================================================================
\section*{Appendix B\,---\,CV(s) (required, 1 page per collaborator)}

% TODO: add one-page CV per collaborator.

\section*{Appendix C\,---\,Citations}

% TODO: add the full references for citations [1]--[3].

\end{document}
